
% eapsthesis class documentation
% Generated documentation for the eapsthesis class
% Location: thesis/eapsthesis-doc.tex
\documentclass[11pt]{article}
\usepackage[T1]{fontenc}
\usepackage[utf8]{inputenc}
\usepackage{geometry}
\geometry{a4paper,margin=2.5cm}
\title{Documentation: eapsthesis class}
\author{Auto-generated from eapsthesis.cls}
\date{2026-08-08}

\frenchspacing
\begin{document}
\maketitle
\tableofcontents
\clearpage

\section{Purpose}
This document describes the options and defaults provided by the eapsthesis class (file: eapsthesis.cls) 
used for BSc/MSc theses in the Department of Earth and Planetary Sciences (D-EAPS), ETH Zurich.
The class itself loads a number of packages that will be available to all documents using the class.
This brings some important constraints with package options that are documented in the corresponding section below.

\section{Loading the class}
Typical use:
\begin{verbatim}
\documentclass[11pt,bscthesis,book]{eapsthesis}
\end{verbatim}

If no thesis type is given, eapsthesis defaults to a BSc thesis in book mode (chapters).

\section{Class options}

\subsection{Supported Options}

The class implements the following explicit options:

\begin{description}
  \item[\texttt{bscthesis}] (default)
    Select Bachelor thesis layout and title page wording. This option is mutually exclusive with \texttt{mscthesis}.
  \item[\texttt{mscthesis}] 
      Select Master's thesis layout and title page wording. When used, the \verb|\specialization| command is 
      displayed on the title page.
      This option is mutually exclusive with \texttt{bscthesis}.
  \item[\texttt{book}] (default: \textbf{true})
    Use a book-like document structure with chapters. When active the class loads the \verb|book| base class 
    with these explicit base-class options: 
    
    \texttt{a4paper, onecolumn, openright, titlepage}.
    The class treats book mode as the two-sided thesis format (defaults to \texttt{twoside} geometry and 
    open-right chapter starts appropriate for printed theses).
  \item[\texttt{report}] (default: false)
    Use a report/article-like structure (sections only). When active the class loads the \texttt{article} 
    base class with these explicit base-class options: \texttt{a4paper, onecolumn, titlepage}.
    The class treats report mode as a \texttt{oneside} (single-sided) format appropriate for section-only documents.
\end{description}

\noindent Additional options from base classes that are supported:

\begin{description}
  \item[\texttt{oneside}] Selects a layout with even left and right margins. The total text area size is the
    same as with \texttt{twoside}, and thus should meet the requirements stated on BSc thesis website.
    This option is automatically selected with \texttt{report}, but can be overridden by explicitly passing \texttt{twoside}.
  \item[\texttt{twoside}] Selects a layour suited for printing and binding. Left and right margins are different,
    as specified on the BSc thesis website. This option is the default for \texttt{book}, but can be overridden
    by explicitly passing \texttt{oneside}.
\end{description}

\subsection{Options explicitly disabled}

The class declares the following base-class options as not allowed (they will be ignored and produce an \texttt{OptionNotUsed}):
\begin{description}
  \item[a5paper, letterpaper, b5paper, legalpaper, executivepaper] Paper format is prescribed.
  \item[twocolumn] Single column is prescribed. 
  \item[openany] This option is not needed. It's the standard behavior for \texttt{report}, and \texttt{book} defaults
    to it when using \texttt{onside}, which is the only case in which this makes sense.
\end{description}

\subsection{Behaviour differences when choosing \texttt{book} vs. \texttt{report}}
The class implements separate frontmatter/mainmatter/backmatter helpers for the two modes. Important differences:
\begin{itemize}
  \item book (chapters): \verb|\frontmatter|, \verb|\mainmatter| and chapter-based table-of-contents and lists are set up. The class loads \texttt{book} with \texttt{openright} and configures headers/footers appropriate for two-sided thesis printing. Page numbering in the frontmatter uses lowercase roman numerals and chapters are used for major ToC entries. Book mode is treated as the default two-sided (\texttt{twoside}) thesis layout.
  \item report (sections-only): \verb|\frontmatter| and \verb|\mainmatter| are defined for section-based documents. TOC and lists are added as sections (not chapters) in the frontmatter and the class loads the \texttt{article} base class (with \texttt{a4paper,onecolumn,titlepage}). Report mode is treated as a oneside (single-sided) layout.
\end{itemize}

Because the class enforces its own base-class options at load time (for example \texttt{a4paper, onecolumn, titlepage} 
are always specified), some package or class options that would change the base-class defaults are reverted to
the eapsthesis defaults above when either \texttt{book} or \texttt{report} is selected.

\subsection{Page geometry and twoside}
The class sets geometry explicitly. If the document is twoside (the \LaTeX\ primitive 
\verb|\ifthenelse{\boolean{@twoside}}| is true), geometry is set to inner=3cm,outer=2cm,top=3cm,bottom=3cm; 
otherwise symmetric margins 2.5cm are used. In typical use: \texttt{book} mode implies two-sided margins 
whereas \texttt{report} implies one-sided margins.

\section{User-facing commands}

The class provides commands to fill in the title page:
\begin{description}
  \item[\texttt{\textbackslash studentid{...}}] Student ID printed on title page.
  \item[\texttt{\textbackslash mainadvisor{<name>}{<affiliation>}}] Main supervisor and affiliation.
  \item[\texttt{\textbackslash coadvisor{<name>}{<affiliation>}}] Additional supervisor and affiliation.
  \item[\texttt{\textbackslash specialization{...}}] Master's specialization text used for MSc title page.
\end{description}

The layout of the title page is fairly rigidly prescribed and shaped by the class accordingly.


\section{Packages loaded by the class}
The class explicitly loads the following packages:
\begin{itemize}
  \item \verb|graphicx| --- \verb|\includegraphics| and friends.
  \item \verb|geometry| --- for page geometry settings (loaded with option \texttt{a4paper}). Settings can be changed
    in the main document.
  \item \verb|fancyhdr| --- for custom headers and footers. Options can be changed in the main document.
  \item \verb|biblatex| --- loaded with options \verb|style=apa, backend=biber| for bibliography handling. See below.
  \item \verb|xifthen| --- provides boolean conditionals and \verb|\newboolean|, \verb|\ifthenelse| used
   by the class. You will likely never need this.
\end{itemize}

That the class loads \texttt{biblatex} brings some important limitations. Most importantly, the citations style
is an option that must be provided a class load time, and cannot be provided later. That means the class
must set that option. At the moment, this is hardcoded to \texttt{apac}, which is appropriate for 
D-EAPS theses. 

\section{Notes and recommendations}
\begin{itemize}
  \item Default mode is equivalent to: \verb|\documentclass[bscthesis,book]{eapsthesis}|.
  \item To produce a Master's thesis use \verb|mscthesis| and set \verb|\specialization{...}| and advisor commands accordingly.
  \item Avoid passing paper-size or column/openside options; the class intentionally restricts or overrides these choices to ensure consistent thesis formatting.
  \item If unusual base-class options are required, review the class file (eapsthesis.cls) and adjust there rather than relying on passing options to \verb|\documentclass|.
\end{itemize}

\section{Implementation notes}
Relevant lines and constructs in \texttt{eapsthesis.cls} (for maintainers): \verb|\DeclareOption{...}|, 
the \verb|\setboolean| defaults at the top of the file, the \verb|\DeclareOption*| forwarding rule, and 
the final \verb|\LoadClass[...]{book}| / \verb|\LoadClass[...]{article}| calls.

Unrecognized options are forwarded to the underlying base class (the class attempts to pass unknown options to 
\texttt{book}). However, eapsthesis later forces specific base-class options when loading the base class 
(see previous section). As a consequence, if you pass options that affect page layout or division behavior 
(paper size, one/two column, openright/openany, titlepage, twoside, etc.), the eapsthesis class may override or 
ignore them in order to enforce its thesis defaults.

\vspace{1em}
\noindent Generated from class file: \texttt{thesis/eapsthesis.cls}.

\end{document}
